\subsection{Interpretation and caveats}
\label{subsec:interpretation}

The three predictions divide the four metrics cleanly. Distinctiveness
and entropy carry no independent growth signal once MA is controlled
for. Both move with MA in Section~\ref{sec:results} and stop mattering
on the growth side, which is what new trade theory predicts. Total
diversity carries a marginally positive growth signal, consistent with
the Romer-Jacobs view that breadth of knowledge drives growth. The two
Mokyr-relevant metrics behave the way the Mokyr channel predicts.
Share-of-useful predicts growth on the extensive margin, driven by
vernacular composition shifts. Within-useful diversity predicts growth
on the intensive margin, driven by Latin scholarly depth. The
cross-language asymmetry is not noise. It is what we would expect if
the channel works through different print media for different
audiences.

Our findings refine the empirical print-and-growth literature on both
margins. \citet{Dittmar2011_information_technology} and
\citet{BoernerRubinSevergnini2021_eer} have shown that whether a city
had a printing press predicts subsequent population growth between
1500 and 1600. Our panel adds the compositional margin. What the
press produced, particularly the share and depth of useful-knowledge
content, predicts growth above and beyond whether a press existed.
Our share-of-useful finding is the city-level, six-language,
1450--1650 analogue of the volume-level, England-published, 1500--1900
corpus-textual finding of \citet{AlmelhemIyigunKennedyRubin2026_qje}.
In their LDA-based analysis of 264{,}443 volumes published in England
in the HathiTrust Digital Library, the volumes that became most
progress-oriented during the Enlightenment were those using language
at the nexus of science and political economy, with
industrial-flavoured works at that same nexus carrying the strongest
signal. The Mokyr channel finds further empirical support in
\citet{SquicciariniVoigtlander2015_qje}, who use subscriber lists for
Diderot's \emph{Encyclop\'edie} as an upper-tail human capital proxy
and show that they predict French city industrialisation a century
later. Their proxy operates on the demand side, ours on the supply
side. The timing of our growth signal also aligns with the
productivity-growth chronology of
\citet{BouscasseNakamuraSteinsson2025_qje}, who estimate that English
productivity growth accelerated around 1600, bracketed by our 1450 and
1650 panel endpoints.

The magnitudes are small. A one-standard-deviation increase in
share-of-useful is associated with about a 1.5 log-point increase in
50-year city growth.

The growth channel may also operate at a longer horizon than the
1450--1650 window suggests. Our reduced-form panel cannot distinguish
equilibrium-conditioned composition coefficients from causal ones, and
the Krugman cumulative-causation dynamics in
\citet{Krugman1991_geography_jpe} explicitly play out over centuries
rather than single 50-year windows. A positive 50-year correlation
could therefore overstate or understate the structural channel.

The Thirty Years' War (1618--1648) shapes the 1600--1650 growth window
for German-territory cities. Appendix~\ref{app:war} reports the
sensitivity, and both Section~\ref{sec:results} composition signals
and Section~\ref{sec:growth} diversity-on-growth strengthen once
war-affected observations are removed.

We do not claim a single channel explains every coefficient. Italian
print shows a significantly negative growth coefficient on diversity
($-0.251^{**}$, Table~\ref{tab:section9_main}) that runs against the
pooled positive, plausibly reflecting the outsized weight of Venezia
in the Italian-print panel (35\% of all Italian titles, see
Section~\ref{sec:geography}) over a period in which Venezia's own
economic position weakened.

\begin{figure}[H]
  \centering
  \includegraphics[width=\textwidth]{../Figures/fig_growth_coefs.pdf}
  \caption{Pooled Barro long-difference coefficients of
    $\Delta_{50}\log(\pop)$ on each composition metric, conditional on
    $\log(\MA)$ and $\log(\pop)$, under the full sample and the
    war-restricted sample. 95\% confidence intervals; dashed vertical
    line at zero. Full regression table in Appendix
    Table~\ref{tab:section9_main}.}
  \label{fig:growth_coefs}
